\section*{Part V: The Market Changes the Meaning of Money}

\subsection*{32. FMCG: Small Differences, Repeated Purchases, Distribution-Driven Choice}

When a person chooses a ₹40 snack, the relevant decision may take seconds. Reference prices, pack size, retailer placement, availability, taste, habit, and visual recognition can dominate detailed reasoning about value. An elaborate explanation does not help a product absent from the shelf. The meaningful test is often paid trial followed by repeat purchase at regular prices, with a full accounting of retailer margins and promotional subsidies.

WTP may vary sharply by occasion: a commuter in a hurry, a parent choosing a weekly staple, and a gift buyer evaluate the same category differently. A 1-time premium paid during a promotion does not prove durable pricing power. The founder must observe repeat behavior under realistic availability and cost.

\subsection*{33. Commodities: Standard Output Does Not Imply Identical Transactions}

For a standardized industrial input, the buyer often knows the market reference price and can switch readily between acceptable suppliers. WTP above that benchmark may attach to timely delivery, traceability, low minimum order, reliability, or credit terms: not to the basic chemical or metal grade. A founder who cannot create a meaningful difference may need a cost advantage rather than a premium narrative.

The most relevant analysis is full landed cost and risk: purchase price, transport, defects, inventory, downtime, financing, and service. In a genuinely homogeneous exchange with many competitive suppliers, sustained markups are limited. The firm's opportunity may be operational excellence rather than increased customer valuation of the underlying object.

\subsection*{34. SaaS and AI Tools: Usage, Ownership, and Risk Diverge}

Software can produce enormous value relative to its distribution cost, but buyers compare workflow alternatives, not source-code complexity. A tool that saves 1 hour a month competes against remaining manual; 1 that prevents expensive errors in an audited process may command far more. Deployment burden, integrations, data governance, reliability, and buyer confidence often dominate novelty.

In AI products, successful adoption can increase compute and inference expense, and quality may be probabilistic. Customers may fear variable bills and unreliable answers. WTP may support a base fee, usage metric, or per-successful-task price, but that is a question for Chapter 16; the prior question is \emph{which output actually matters and what alternative cost it replaces}.

\subsection*{35. Professional Services: Expertise, Accountability, and Career Safety}

The value of a specialist's 6 hours cannot be inferred from a junior consultant's 6 hours. A diagnostic that avoids a large loss may be worth far more than the labor required to produce it. Still, the buyer must believe the diagnosis, be able to use it, and know what responsibilities the provider assumes. Time saved in a proposal is not savings realized in operations.

High WTP for professional services often depends on a narrow task, scarce judgment, reputational exposure, and the ability to make an ambiguous problem governable. A vague ``strategic transformation'' proposal may be harder to justify than a bounded project with explicit evidence and acceptance conditions.

\subsection*{36. Marketplaces: WTP Is Interdependent Across Sides}

The buyer may value a marketplace because enough suitable sellers participate; sellers may value it because buyers are present and transactions clear reliably. These are interdependent valuations. A founder cannot optimize 1 side's immediate payment in isolation if doing so causes the other side to disappear.

Analyze each side's alternative, search cost, trust cost, and ability to transact elsewhere. The platform may subsidize 1 side to create liquidity, but only if the subsidy has a credible path to system value. A database of registered users does not produce willingness to pay when reliable matches fail to occur.

\subsection*{37. Regulated and High-Stakes Purchases: Willingness Is Bounded by Legitimacy}

Financial, medical, legal, and other high-stakes services face obligations that cannot be answered by ``the customer would pay.'' Buyers may have impaired information, urgent vulnerability, or limited ability to evaluate expertise. Reimbursement systems, professional ethics, and regulation also shape the feasible transaction. The firm must pursue legitimate value under clear standards, not optimize extraction from fear.

A founder seeking a durable professional institution benefits from transparent boundaries: what outcome is controlled, what uncertainty remains, what evidence supports the intervention, who is qualified to provide it, and what happens when the service fails.
